\documentclass[10pt,letterpaper]{article}
\usepackage[ficha]{../latex_common/estilo_monitoria}

\fancyhead[L]{\footnotesize\sffamily\color{GrisTexto}\textbf{Ficha Resumen 02 --- Traductor DFD $\leftrightarrow$ PSeInt $\leftrightarrow$ C++}}
\fancyhead[R]{\footnotesize\sffamily\color{Turquesa}\textbf{Monitoría Algoritmos 30h}}

\begin{document}
\thispagestyle{empty}

\begin{tcolorbox}[enhanced,colback=AzulMarino,colframe=AzulMarino,arc=4pt,top=3.5mm,bottom=3.5mm,left=4mm,right=4mm,boxrule=0pt]
  \tcbox[on line,colback=Turquesa,colframe=Turquesa,arc=2pt,boxsep=1pt,left=4pt,right=4pt]{\footnotesize\sffamily\bfseries\color{white} FICHA IMPRIMIBLE 02}\hspace{6pt}%
  {\large\bfseries\sffamily\color{white} Traductor Maestro: DFD $\leftrightarrow$ PSeInt $\leftrightarrow$ C++ Básico}\\[2pt]
  {\footnotesize\sffamily\color{AzulSuave} El puente exacto de instrucciones, tipos de datos, ciclos y trampas de traducción}
\end{tcolorbox}

\begin{center}
  \includegraphics[width=0.58\textwidth,height=4.2cm,keepaspectratio]{ilus_03_puente_dfd_pseint_cpp.jpg}\hfill
  \includegraphics[width=0.40\textwidth,height=4.2cm,keepaspectratio]{mem_02_puente_ciclos.png}
\end{center}

\vspace{-6pt}
\begin{center}
  \scriptsize
  \renewcommand{\arraystretch}{1.18}
  \rowcolors{2}{GrisSuave}{white}
  \begin{tabularx}{\textwidth}{>{\bfseries\sffamily\color{AzulMarino}}p{2.7cm} >{\raggedright\arraybackslash}p{2.8cm} >{\ttfamily\raggedright\arraybackslash}X >{\ttfamily\raggedright\arraybackslash}X}
    \rowcolor{AzulMarino}
    \textcolor{white}{\textbf{Concepto}} & \textcolor{white}{\textbf{DFD}} & \textcolor{white}{\textbf{\sffamily PSeInt (español)}} & \textcolor{white}{\textbf{\sffamily C++ básico (C++17)}} \\
    \midrule
    Estructura base & Óvalos Inicio / Fin & Algoritmo Nom ... FinAlgoritmo & int main() \{ ... return 0; \} \\
    Tipos de datos & Implícitos & Entero, Real, Caracter, Logico, Cadena & int, double, char, bool, string \\
    Constante & Valor fijo & (No nativa / convención MAYÚSC) & const int MAX = 100; \\
    Entrada teclado & Paralelogramo / Trapecio & Leer x; & cin >> x; \textit{(texto: getline)} \\
    Salida pantalla & Cono / Pantalla & Escribir "Total: ", t; & cout << "Total: " << t << endl; \\
    Asignación & Rectángulo & x <- a + b; & x = a + b; \\
    Operadores lógicos & Y / O / NO & Y, O, NO, MOD, =, <> & \&\&, ||, !, \%, ==, != \\
    Si / SiNo & Rombo (Sí / No) & Si c Entonces ... SiNo ... FinSi & if (c) \{ ... \} else \{ ... \} \\
    Selección múltiple & Rombos en cascada & Segun op Hacer 1: ... FinSegun & switch(op) \{ case 1: ... break; \} \\
    Mientras & Ciclo entrada arriba & Mientras c Hacer ... FinMientras & while (c) \{ ... \} \\
    Repetir vs Do-While & Condición al final & Repetir ... Hasta Que (termina V) & do \{ ... \} while (sigue V); \\
    Para & Contador automático & Para i <- 0 Hasta n-1 Hacer & for (int i = 0; i < n; i++) \\
    Vector estático & Arreglo fijo & Dimension v[100]; (base 1 ó 0) & const int MAX=100; int v[MAX]; \\
    Matriz estática & Tabla 2D fija & Dimension m[10,10]; & int m[10][10]; \\
    Subprograma & Subrutina & SubProceso f(x Por Referencia) & void f(int \&x) \{ ... \} \\
    \bottomrule
  \end{tabularx}
\end{center}

\vspace{-6pt}
\begin{peligro}{Las 5 Trampas Mortales al Traducir de PSeInt a C++}
\scriptsize
\begin{enumerate}[leftmargin=12pt,itemsep=1.5pt]
  \item \textbf{Ciclo invertido:} En PSeInt, \codigo{Repetir ... Hasta Que op = 0} \textbf{TERMINA} cuando \codigo{op == 0}. En C++, \codigo{do \{ ... \} while (op != 0);} \textbf{CONTINÚA} mientras sea verdadero. \textbf{¡Niega la condición al traducir!}
  \item \textbf{Índices 1 vs 0:} PSeInt por defecto indexa de \codigo{1} a \codigo{n}. En C++ siempre va de \codigo{0} a \codigo{n - 1}.
  \item \textbf{División entera:} En PSeInt \codigo{5 / 2} da \codigo{2.5}; en C++ \codigo{5 / 2} da \codigo{2}. Usa \codigo{5.0 / 2} o \codigo{static\_cast<double>(a) / b}.
  \item \textbf{Rangos encadenados:} Nunca escribas \codigo{10 <= x <= 20} en C++; escribe \codigo{10 <= x \&\& x <= 20}.
  \item \textbf{Igualdad:} En PSeInt comparas con \codigo{=}; en C++ comparar es \codigo{==} (un solo \codigo{=} asigna).
\end{enumerate}
\end{peligro}

\end{document}
